\subsection{A scaling non-compact QCD axion --- arXiv:2606.17019}

\textbf{Authors:} Georgios K. Karananas, Mikhail Shaposhnikov

\paragraph{主な主張}
非コンパクトなQCD axionに指数ポテンシャルを加え、QCD期より十分前に放射と同じscaling解へ到達させることで、インフレーション起源の等曲率摂動を動的に減衰させる機構を提案した~\cite{Karananas:2026scaling}。
さらに、QCDポテンシャルへ到達する速度に応じた暗黒物質生成と、指数項が残すstrong CP位相を評価した。

\paragraph{新規性}
既知の指数ポテンシャルのscaling解~\cite{Copeland:1997scaling}そのものではなく、同じ指数項によって等曲率摂動の抑制、kinetic misalignmentへつながる速度、残留CP violationを一つのQCD axion模型に結びつけた点にある。

\paragraph{位置付け}
scaling機構を暗黒物質・等曲率制約・電気双極子モーメント探索に結びつけた一場模型の例であり、非周期的なポテンシャルを前提とするため、周期的axionを含む$SL(2,\mathbb{Z})$二場模型の角方向の収束を直接保証するものではない。
